% !PRINTED 34
\weberpage{34}
%<<CONT
\weq{3}{1 = A_1^{\alpha_1' - \alpha_1} A_2^{\alpha_2' - \alpha_2} \ldots A_n^{\alpha_n' - \alpha_n},}
thus a representation of the element 1. Hence by (2) we may set:
\[
\alpha_1' = \alpha_1 + h_1,\ \alpha_2' = \alpha_2 + h_2 \ldots \alpha_n' = \alpha_n + h_n.
\]
Consequently there can also be no more distinct representations of the element $\Theta$ than the number of representations of the element 1 amounts to. We denote the number of these representations by $h$. In all, however, the number of the various possible determinations of the $\alpha$ in (1) is equal to the product $a_1 a_2 \ldots a_n$, and since $n$ is the number of elements of $S$, it follows that
\weq{4}{n h = a_1 a_2 \ldots a_n.}
From the formula (4) results the theorem:

\gesperrt{2. If $r$ is a prime dividing the degree $n$ of the group, then there is in $S$ an element of degree $r$.}

For from (4) one sees first of all that one of the degrees $a_1, a_2, \ldots a_n$ is divisible by $r$. If then, say, $a_1 = rk$, then $A_1^k$ is an element of degree $r$.

\gesperrt{3. If $A, B \ldots$ are any elements whatever in $S$ of the degrees $a, b \ldots$, then $AB \ldots$ is also an element in $S$, and the degree of this product is a divisor of the least common multiple $m$ of $a, b \ldots$}

For we have
\[
(AB \ldots)^m = A^m B^m \ldots = 1,
\]
and consequently $m$ is a multiple of the degree of $AB \ldots$

\gesperrt{4. If the degree $n$ of the group $S$ is decomposed into two factors $n = ab$ such that $a$ and $b$ are relatively prime, then there are in $S$ exactly $a$ elements $A$ whose degree is a divisor of $a$, and $b$ elements $B$ whose degree is a divisor of $b$, and in the form}
\weq{5}{\Theta = AB}
\gesperrt{all the elements of $S$ are contained, and each only once.}

To see the correctness of this theorem, we make the following consideration. The totality $\mathfrak{A}$ of the elements $A$ whose
%CONT>>
